\chapter[Transacciones y Control de Concurrencia]{\emj{🔒}~Transacciones y Control de Concurrencia}

\section[ACID]{\emj{⚗️}~ACID}

Una transacción es una secuencia de operaciones que el DBMS trata como una unidad indivisible: o todas se completan exitosamente, o ninguna tiene efecto.
Las propiedades ACID son las garantías que hacen a las transacciones confiables incluso en presencia de fallos del sistema, cortes de energía o acceso concurrente de miles de usuarios simultáneos.

ACID no es magia: es el resultado de mecanismos concretos implementados por el DBMS.
Atomicidad usa undo logs; Durabilidad usa WAL (Write-Ahead Log); Aislamiento usa locks o MVCC; Consistencia es responsabilidad compartida entre el DBMS (constraints) y la aplicación (lógica de negocio).

\begin{teoriabox}{ACID explicado con transferencia bancaria}
Transferir \$1{,}000 de la cuenta A a la cuenta B requiere dos operaciones:
(1) restar \$1{,}000 de A y (2) sumar \$1{,}000 a B.

\begin{itemize}
  \item \textbf{Atomicidad:} si el servidor se cae entre (1) y (2), al reiniciar el DBMS deshace la operación (1) usando el undo log. El dinero no desaparece ``a la mitad''. O ambas operaciones ocurren, o ninguna.

  \item \textbf{Consistencia:} la suma total del dinero en el sistema es idéntica antes y después. Ningún constraint se viola durante ni después de la transacción. Si A no tiene saldo suficiente y hay un CHECK constraint, la transacción falla completa.

  \item \textbf{Aislamiento:} si otra transacción consulta el saldo de A mientras la transferencia está en curso, no ve el estado intermedio (A descontado, B sin acreditar). Ve o el estado antes de la transferencia, o el estado después — nunca el estado a medias.

  \item \textbf{Durabilidad:} una vez que el \texttt{COMMIT} regresa con éxito, los cambios están en el WAL en disco. Aunque el servidor se apague inmediatamente, al reiniciar el DBMS replaya el WAL y los datos están ahí.
\end{itemize}
\end{teoriabox}

\section[Niveles de aislamiento]{\emj{🛡️}~Niveles de aislamiento}

El aislamiento perfecto (serializable) implicaría que las transacciones se ejecuten como si fueran estrictamente secuenciales.
Esto sería correcto pero extremadamente lento: no habría concurrencia real.
Los niveles de aislamiento son el compromiso entre correctitud y rendimiento: mayor aislamiento = más correctitud, menos concurrencia.

El estándar SQL define cuatro anomalías que los niveles de aislamiento previenen o permiten:
\begin{itemize}
  \item \textbf{Dirty Read:} una transacción lee datos de otra que aún no hizo COMMIT. Si la segunda hace ROLLBACK, la primera leyó datos que nunca existieron.
  \item \textbf{Non-Repeatable Read:} una transacción lee la misma fila dos veces y obtiene valores distintos porque otra transacción la modificó entre lecturas.
  \item \textbf{Phantom Read:} una transacción ejecuta la misma query dos veces y obtiene filas distintas porque otra transacción insertó o eliminó filas que cumplen el filtro.
  \item \textbf{Serialization Anomaly:} el resultado de ejecutar transacciones concurrentemente no es posible obtenerlo con ningún orden serial de las mismas transacciones.
\end{itemize}

Los cuatro niveles de aislamiento en PostgreSQL:
\begin{itemize}
  \item \textbf{Read Uncommitted:} en PostgreSQL se trata igual que Read Committed (no existe dirty read en PG).
  \item \textbf{Read Committed (default PG):} cada sentencia ve solo los datos que estaban confirmados al momento de iniciar esa sentencia. Permite non-repeatable reads.
  \item \textbf{Repeatable Read:} la transacción ve un snapshot fijo del momento en que inició. Mismo valor al releer la misma fila. Previene dirty y non-repeatable reads. En PG también previene phantom reads.
  \item \textbf{Serializable:} garantía máxima. PG detecta ciclos de dependencia entre transacciones concurrentes y aborta las que crearían anomalías de serialización.
\end{itemize}

\begin{tipbox}{¿Qué nivel usar en la práctica?}
\textbf{Read Committed} (default) cubre el 95\% de los casos de aplicaciones web.
Usa \textbf{Repeatable Read} cuando una transacción hace múltiples lecturas del mismo dato y la inconsistencia entre ellas causaría lógica incorrecta (reportes financieros, cálculos que dependen de varios SELECTs).
Usa \textbf{Serializable} solo cuando la correctitud lógica de la operación dependa de ausencia total de interferencia concurrente. Tiene overhead de detección de conflictos y puede generar más errores que requieren retry.
\end{tipbox}

\section[Locks y Deadlocks]{\emj{⚠️}~Locks y Deadlocks}

PostgreSQL usa \textbf{MVCC} (Multi-Version Concurrency Control) como mecanismo principal de concurrencia.
En MVCC, en lugar de bloquear filas para lectura, el motor mantiene múltiples versiones de cada fila.
Cada transacción ve un snapshot consistente de la BD al momento de su inicio, sin necesidad de locks de lectura.
Esto elimina el clásico problema de lecturas que bloquean escrituras y viceversa.

Sin embargo, dos escrituras concurrentes sobre la misma fila sí requieren coordinación.
PostgreSQL usa locks a nivel de fila para escrituras: la segunda escritura espera a que la primera haga COMMIT o ROLLBACK.

Un \textbf{deadlock} ocurre cuando dos o más transacciones forman un ciclo de espera: T1 espera un lock que tiene T2, y T2 espera un lock que tiene T1.
Ninguna puede avanzar indefinidamente.
PostgreSQL detecta estos ciclos periódicamente y aborta la transacción más joven (menor costo de rollback), retornando un error al cliente para que reintente.

\begin{warningbox}{Cómo prevenir deadlocks}
La causa raíz de casi todos los deadlocks es acceder a recursos (filas, tablas) en órdenes distintos en transacciones diferentes.
Solución: establecer y documentar un \textbf{orden canónico} de acceso a recursos en toda la aplicación.
Si T1 y T2 siempre adquieren locks en el orden Recurso A $\to$ Recurso B, el deadlock es matemáticamente imposible.
Además: mantener las transacciones lo más cortas posible y nunca esperar input del usuario dentro de una transacción abierta.
\end{warningbox}

\begin{lstlisting}[caption={BEGIN, COMMIT, ROLLBACK, SAVEPOINT}]
BEGIN;

UPDATE cuentas SET saldo = saldo - 1000 WHERE id = 1;
UPDATE cuentas SET saldo = saldo + 1000 WHERE id = 2;

-- Si algo falla:
ROLLBACK;
-- Si todo OK:
COMMIT;

-- SAVEPOINT: rollback parcial
BEGIN;
UPDATE inventario SET stock = stock - 5 WHERE producto_id = 10;
SAVEPOINT antes_descuento;
UPDATE precios SET precio = precio * 0.9 WHERE producto_id = 10;
ROLLBACK TO SAVEPOINT antes_descuento;
COMMIT;
\end{lstlisting}

\begin{lstlisting}[caption={SELECT FOR UPDATE y nivel de aislamiento}]
-- Bloquear fila para actualización (evita race condition)
BEGIN;
SELECT saldo FROM cuentas WHERE id = 1 FOR UPDATE;
UPDATE cuentas SET saldo = saldo - 500 WHERE id = 1;
COMMIT;

-- Cambiar nivel de aislamiento
BEGIN TRANSACTION ISOLATION LEVEL REPEATABLE READ;
SELECT total FROM reporte WHERE mes = '2024-01';
COMMIT;
\end{lstlisting}

\begin{lstlisting}[caption={Detectar y matar conexiones colgadas}]
-- Ver transacciones activas hace más de 5 minutos
SELECT pid, usename, state, query_start,
       NOW() - query_start AS duracion,
       LEFT(query, 80) AS query_corta
FROM   pg_stat_activity
WHERE  state != 'idle'
  AND  query_start < NOW() - INTERVAL '5 minutes';

-- Terminar transacción colgada
SELECT pg_terminate_backend(12345);
\end{lstlisting}

\begin{entrevistabox}
\begin{enumerate}
  \item Explica ACID con un ejemplo de transferencia bancaria.
  \item ¿Qué es un deadlock y cómo lo detecta/resuelve PostgreSQL?
  \item ¿Cuál es la diferencia entre Read Committed y Repeatable Read?
\end{enumerate}
\end{entrevistabox}
